\documentclass[preview]{standalone}

\usepackage{amsmath}
\usepackage{amssymb}
\usepackage{stellar}
\usepackage{definitions}

\begin{document}

\id{psicologia-organizzazione-percettiva}
\genpage

\section{Organizzazione percettiva}

\begin{snippetdefinition}{processi-organizzazione-percettiva-definition}{Processi di organizzazione percettiva}
    I \emph{processi di organizzazione percettiva} sono processi di integrazione
    delle informazioni sensoriali che permettono di costruire una percezione
    coerente a partire dall'attività dei recettori.
\end{snippetdefinition}

\begin{snippet}{gestalt-organizzazione-percettiva}
    La teoria della Gestalt ha descritto la percezione come un'organizzazione
    regolata da leggi. Distinzione figura-sfondo, percezione di profondità e
    movimento, integrazione spaziale e integrazione temporale dipendono da
    principi percettivi di base.
\end{snippet}

\begin{snippetdefinition}{articolazione-figura-sfondo-definition}{Articolazione figura-sfondo}
    L'\emph{articolazione figura-sfondo} è il processo percettivo attraverso
    cui, all'interno di una scena, un elemento viene percepito come figura e il
    resto come sfondo. La figura ha forma ed estensione definite, mentre lo
    sfondo è meno differenziato e continua dietro la figura.
\end{snippetdefinition}

\begin{snippetdefinition}{figure-reversibili-definition}{Figure reversibili}
    Le \emph{figure reversibili} sono configurazioni percettive instabili e
    ambigue nelle quali figura e sfondo possono invertirsi e alternarsi.
\end{snippetdefinition}

\begin{snippet}{principi-raggruppamento-gestalt}
    Max Wertheimer formulò diversi principi di raggruppamento percettivo:
    \begin{itemize}
        \item \emph{vicinanza}: a parità di condizioni, si unificano gli
        elementi vicini;
        \item \emph{somiglianza}: a parità di condizioni, si unificano gli
        elementi simili;
        \item \emph{buona direzione}: a parità di condizioni, si unificano gli
        elementi che presentano continuità di direzione;
        \item \emph{chiusura}: a parità di condizioni, si unificano gli
        elementi che tendono a chiudersi tra loro;
        \item \emph{destino comune}: a parità di condizioni, si unificano gli
        elementi che condividono la stessa direzione di movimento.
    \end{itemize}
\end{snippet}

\begin{snippetdefinition}{integrazione-spaziale-definition}{Integrazione spaziale}
    L'\emph{integrazione spaziale} è il processo attraverso cui il sistema
    percettivo combina informazioni provenienti da differenti collocazioni
    spaziali.
\end{snippetdefinition}

\begin{snippetdefinition}{integrazione-temporale-definition}{Integrazione temporale}
    L'\emph{integrazione temporale} è il processo attraverso cui il sistema
    percettivo combina informazioni raccolte in differenti momenti nel tempo.
\end{snippetdefinition}

\begin{snippetdefinition}{estensione-confini-definition}{Estensione dei confini}
    L'\emph{estensione dei confini} è un errore percettivo e mnestico per cui,
    dopo aver osservato una scena, le persone tendono a ricordare di aver visto
    una porzione più ampia rispetto a quella effettivamente presentata.
\end{snippetdefinition}

\begin{snippetexample}{estensione-confini-example}{Estensione dei confini}
    Quando osserviamo una fotografia in primo piano, possiamo ricordarla come
    se mostrasse una porzione più ampia della scena. La conoscenza del mondo
    viene usata per estendere mentalmente i confini dell'inquadratura oltre ciò
    che era direttamente visibile.
\end{snippetexample}

\includesnpt[width=70\%,src=/snippet/static-psychology/boundary-extension-example.jpg]{centered-img}

\begin{snippetdefinition}{cecita-cambiamento-definition}{Cecità al cambiamento}
    La \emph{cecità al cambiamento}, o \emph{change blindness}, è la difficoltà
    a notare cambiamenti tra diverse prospettive o versioni di una stessa
    scena.
\end{snippetdefinition}

\begin{snippetdefinition}{percezione-movimento-definition}{Percezione del movimento}
    La \emph{percezione del movimento} è la capacità di combinare informazioni
    provenienti da osservazioni visive successive per percepire lo spostamento
    di uno stimolo nello spazio.
\end{snippetdefinition}

\begin{snippetexample}{fenomeno-phi-example}{Fenomeno phi}
    Il \emph{fenomeno phi} si verifica quando due punti luce statici, collocati
    in posizioni differenti, vengono accesi e spenti in modo alternato a un
    ritmo adeguato. L'osservatore non percepisce due punti separati, ma un unico
    movimento apparente lungo un percorso semplice.
\end{snippetexample}

\begin{snippetdefinition}{percezione-profondita-definition}{Percezione della profondità}
    La \emph{percezione della profondità} è la capacità di percepire la distanza
    e la posizione degli oggetti nello spazio tridimensionale.
\end{snippetdefinition}

\begin{snippetdefinition}{indizi-binoculari-profondita-definition}{Indizi binoculari di profondità}
    Gli \emph{indizi binoculari di profondità} sono informazioni sulla
    profondità che derivano dal confronto tra le informazioni visive fornite dai
    due occhi.
\end{snippetdefinition}

\begin{snippetdefinition}{disparita-retinica-definition}{Disparità retinica}
    La \emph{disparità retinica} è lo scarto tra le posizioni orizzontali
    retiniche di immagini corrispondenti nei due occhi. Fornisce informazioni
    sulla profondità perché il grado di disparità dipende dalla distanza
    relativa tra l'oggetto e l'osservatore.
\end{snippetdefinition}

\begin{snippetdefinition}{convergenza-definition}{Convergenza}
    La \emph{convergenza} è un'informazione binoculare di profondità basata sul
    movimento degli occhi verso l'interno quando si osserva uno stimolo vicino.
\end{snippetdefinition}

\begin{snippetdefinition}{movimento-parallasse-definition}{Movimento di parallasse}
    Il \emph{movimento di parallasse} è un'informazione di profondità
    determinata dalla velocità e dalla direzione del movimento relativo degli
    oggetti nell'immagine retinica.
\end{snippetdefinition}

\begin{snippetexample}{movimento-parallasse-example}{Movimento di parallasse}
    Quando si viaggia in automobile, gli oggetti vicini sembrano muoversi più
    rapidamente rispetto agli oggetti lontani, che appaiono più statici. Questa
    differenza di movimento relativo contribuisce alla percezione della
    profondità.
\end{snippetexample}

\begin{snippetdefinition}{indizi-monoculari-profondita-definition}{Indizi monoculari di profondità}
    Gli \emph{indizi monoculari di profondità} sono informazioni sulla
    profondità che possono essere ricavate anche usando un solo occhio.
\end{snippetdefinition}

\begin{snippet}{principali-indizi-monoculari}
    I principali indizi monoculari di profondità includono:
    \begin{itemize}
        \item \emph{occlusione}: quando un oggetto opaco si sovrappone
        parzialmente a un altro, l'oggetto nascosto viene percepito come più
        lontano;
        \item \emph{dimensione relativa}: oggetti della stessa dimensione
        collocati a distanze differenti proiettano immagini retiniche di
        dimensioni differenti;
        \item \emph{prospettiva lineare}: la distanza viene inferita dalla
        convergenza apparente delle linee parallele;
        \item \emph{gradienti di trama}: la densità della trama aumenta quanto
        più una superficie si allontana.
    \end{itemize}
\end{snippet}

\begin{snippetdefinition}{costanza-percettiva-definition}{Costanza percettiva}
    La \emph{costanza percettiva} è il fenomeno per cui le proprietà degli
    oggetti vengono percepite come stabili nonostante le variazioni della
    stimolazione sensoriale.
\end{snippetdefinition}

\begin{snippetdefinition}{costanza-grandezza-definition}{Costanza di grandezza}
    La \emph{costanza di grandezza} è la capacità di percepire le dimensioni
    effettive di un oggetto nonostante le variazioni di grandezza della sua
    immagine retinica.
\end{snippetdefinition}

\begin{snippetdefinition}{costanza-forma-definition}{Costanza di forma}
    La \emph{costanza di forma} è la capacità di percepire la forma di un
    oggetto come stabile anche quando esso è inclinato rispetto all'osservatore
    e produce un'immagine retinica differente.
\end{snippetdefinition}

\begin{snippetdefinition}{costanza-luminosita-definition}{Costanza di luminosità}
    La \emph{costanza di luminosità} è la tendenza a percepire il bianco, il
    grigio o il nero degli oggetti come costanti nonostante le variazioni di
    illuminazione.
\end{snippetdefinition}

\end{document}
